\newpage\handout[true]
{\textsc{Math 2106-B: Foundations of Mathematical Proof}}{Fall 2025}
{Homework 1}
{\textit{Prof.} \href{https://ceheitsch.github.io/webpage/}{Dr. Christine Heitsch}}{Due Date: August 28th}
{Math 2106-B: HW1}
{\large
  \noindent
  \textbf{Name:} \HandoutAuthor \smallskip \\

  \noindent
  \textbf{GTID:} 903743506 \smallskip \\

  \noindent
  %%% List anyone with whom you discussed the problems
  \textbf{Collaborators:} None. This material reflects independent preparation. \smallskip \\

  \noindent
  %%% List all resources used OTHER than the textbook, lecture notes, 
  %%% webwork problems, ICA questions, and previous homework assignments.
  \textbf{Outside resources:} This work was informed by MATH 2106 course materials 
  (ICAs: in-class activities and WebWork contents) from \HandoutInstitution\ (the course textbook is 
  Chartrand \parencite{chartrand_mathematical_2018}).   
  Outside references specifically include
  textbooks written by 
  Grimaldi \parencite{grimaldi_discrete_2004} (used in MATH 3012) and 
  Rosen \parencite{rosen_discrete_2019} (used in CS 2050). \smallskip \\

  \noindent
  \textbf{Notice}: The answers contain cross-references throughout this document as clickable hyperlinks in most PDF viewers. However, the Gradescope PDF view might not support clickable hyperlinks.
}
The following table presents the standard logical equivalences from these sources together with 
the author's own (or the student, myself, Jaehoon Song) clarifications and labels where applicable.


\begin{table}[h!]
  \centering
  \renewcommand{\arraystretch}{1.2}
  \setlength{\tabcolsep}{8pt}
  \begin{tabular}{|p{5.0cm}|p{5.0cm}|p{5.2cm}|}
    \hline
    \rowcolor[HTML]{E0E0E0}
    \multicolumn{2}{|l|}{\textbf{Logical Equivalences}} & \textbf{Name} \\
    \hline
    $\,p \AND q \equiv q \AND p$ & $\,p \OR q \equiv q \OR p$ & Commutative Law \\
    \hline
    $\, (p \AND q) \AND r \equiv p \AND (q \AND r)$ & $\, (p \OR q) \OR r \equiv p \OR (q \OR r)$ & Associative Law \\
    \hline
    $\,p \AND (q \OR r) \equiv (p \AND q) \OR (p \AND r)$ & $\,p \OR (q \AND r) \equiv (p \OR q) \AND (p \OR r)$ & Distributive Law \\
    \hline
    $\,\overline{p \AND q} \equiv \sNOT{p} \OR \sNOT{q}$ & $\,\overline{p \OR q} \equiv \sNOT{p} \AND \sNOT{q}$ & DeMorgan's Law \\
    \hline
    \multicolumn{2}{|l|}{$\,\overline{\overline{p}} \equiv p$} & Double Negation law \\
    \hline
    $\,p \AND \boldsymbol{T} \equiv p$ & $\,p \OR \boldsymbol{F} \equiv p$ & Identity Law \\
    \hline
    $\,p \OR \boldsymbol{T} \equiv \boldsymbol{T}$ & $\,p \AND \boldsymbol{F} \equiv \boldsymbol{F}$ & Domination Law \\
    \hline
    $\,p \AND p \equiv p$ & $\,p \OR p \equiv p$ & Idempotent Law \\
    \hline
    \multicolumn{2}{|l|}{$\,p \THEN q \equiv \sNOT{p} \OR q$} & Implication Equivalence \\
    \hline
    \multicolumn{2}{|l|}{$\,\overline{p \THEN q} \equiv p \AND \sNOT{q}$} & Negation of Implication \\
    \hline
    \multicolumn{2}{|l|}{$\,p \THEN q \equiv \sNOT{q} \THEN \sNOT{p}$} & Contrapositive Equivalence \\
    \hline
    \multicolumn{2}{|l|}{$\,p \OR \sNOT{p} \equiv \boldsymbol{T}$} & Law of Tautology \\
    \hline
    \multicolumn{2}{|l|}{$\,p \AND \sNOT{p} \equiv \boldsymbol{F}$} & Law of Contradiction \\
    \hline
  \end{tabular}
\end{table}\newpage
% The following definitions and axioms are presented from standard mathematical sources (\textbf{Number Theory}) together with the author's own (or the student, myself, Jaehoon Song) clarifications and labels where applicable.

\begin{define}{Integers}[def:integers][red]
  $\setZ$: the set of all integers
\end{define}

\begin{define}{Rational Numbers}[def:rational-numbers][red]
  $\setQ$: the set of all rational numbers
\end{define}

\begin{define}{Real Numbers}[def:real-numbers][red]
  $\setR$: the set of all real numbers
\end{define}

\begin{define}{Natural Numbers}[def:natural-numbers][red]
  $\setN$ or $\setZ_{>0}$: the set of all natural numbers which are the set of positive integers (not including $0$)
\end{define}

\begin{define}{Positive Real Numbers}[def:positive-real-numbers][red]
  $\setR_{>0}$: the set of all positive real numbers
\end{define}

\begin{define}{Non-negative Real Numbers}[def:non-negative-real-numbers][red]
  $\setR_{\geq 0}$: the set of all non-negative real numbers
\end{define}


\newpage

The following definitions and claims extend the understanding of integers.

Here are the ICA2 Definitions and Axioms as follows.

\begin{define}{Even Integer}[def:even-integer][red]
  An integer $n$ is \underbar{even} if there exists 
  (existential quantification logic $\exists$) an integer $k$ such that
  \[
  n = 2k.
  \]
\end{define}

\begin{define}{Odd Integer}[def:odd-integer][red]
  An integer $n$ is \underbar{odd} if 
  there exists some integer $k$ such that
  \[
  n = 2k + 1.
  \]
\end{define}

\begin{define}{Closure under Addition}[def:closure-addition][red]
  The sum of two integers is again an 
  integer. (Closure under addition)
\end{define}

\begin{define}{Closure under Multiplication}[def:closure-multiplication][red]
  The product of two integers is again 
  an integer. (Closure under multiplication)
\end{define}

\begin{define}{Odd if and only if Not Even}[def:odd-iff-not-even][red]
  An integer is odd if and only if it 
  is not even.
\end{define}

Here are the ICA3 Definitions and Axioms as follows.

\begin{define}{Universal Closure under Addition}[def:universal-closure-addition][red]
  The sum of two integers (Universal quantification for all, for every, for each) is again an integer. (Closure under addition)
\end{define}

\begin{define}{Universal Closure under Multiplication}[def:universal-closure-multiplication][red]
  The product of two integers (Universal quantification for all, for every, for each) is again an integer. (Closure under multiplication)
\end{define}

\begin{define}{Even Integer Square}[def:even-integer-square][red]
  An integer $n$ is even if and only if $n^2$ is even.
\end{define}

\newpage

The following claims are given in the reference as well.

\begin{define}{Sum of Even Integers}[claim:sum-even-integers][red]
  The sum of two even integers is again an even integer.
  \begin{subprove}[red]
    Let $x$ be an integer and $y$ be an integer ($\forall x,y \in \setZ$). 
    Suppose $x$ is even and $y$ is even. Then, by \nameref{def:even-integer}, there exists an integer 
    $k$ and an integer $j$ such that $x=2k$ and $y=2j$. Then, 
    $x+y=\left(2k\right)+\left(2j\right)=2\left(k+j\right)$. By \nameref{def:closure-addition}, since $k$ and 
    $j$ are integers, then so is $\left(k+j\right)$. Hence, by definition, $x+y$ is an 
    even integer.
  \end{subprove}
\end{define}

The following definitions extend the understanding of integers including divisibility.

\begin{define}{Divisibility}[def:divisibility][red]
  An integer $n$ divides an integer $m$, denoted $n \LDIV m$, if there 
  exists an integer $k$ such that $m=n k$. Otherwise, $n \nmid m$.
\end{define}

The following definitions extend the classification of real numbers into rational and irrational numbers.

\begin{define}{Rational Number}[def:rational-number][red]
  A real number $x$ is \underbar{rational} if there exist integers $a$ and $b$ with $b \neq 0$ such that $bx = a$.
  Otherwise, $x$ is \underbar{irrational}.

  \textbf{Note}: Compare the rational number definition with the even/odd axiom \inlinequote{odd iff not even} from \nameref{def:odd-iff-not-even}.
\end{define}

The following claims are given in the reference as well.

\begin{define}{Irrationality of $\sqrt{2}$}[claim:sqrt2-irrational][red]
  $\sqrt{2}$ is irrational.
  \begin{subprove}[red]
    Suppose $\sqrt{2}$ is rational. Then, by \nameref{def:rational-number}, there exists $a,b\in \setZ$ 
    with $b\ne 0$ such that $b\cdot\sqrt{2}=a$. We may further assume that any factors 
    common to $a$ and $b$ have already been removed.\\

    Then, squaring both sides, this yields $2\cdot b^2 = a^2$.
    Since products of integers are again integers by \nameref{def:closure-multiplication}, we have that $a^2,b^2\in \setZ$.
    Thus, by definition, $a^2$ is even. Therefore, by \nameref{def:even-integer-square}, $a$ is even. 
    Thus, there exists $c\in\setZ$ such that $a=2c$. Then
    $$
    2b^2=a^2={2c}^2=4c^2
    $$
    Then, $b^2=2c^2$. Then, as before, $b^2$ is even, so $b$ is even.
    However, $a$ and $b$ both being even contradicts the fact that any factors 
    common to $a$ and $b$ have already been removed.
  \end{subprove}
\end{define}






%%%%%%%%%%%%%%%%%%% New style (refactor model)
\begin{define}{Relatively Prime Integers (Coprime)}[def:relatively-prime][red]
  Let $x,y \in \mathbb{Z}$. The integers $x$ and $y$ are said to be
  \textbf{relatively prime} if and only if
  $\gcd(x,y) = 1$.
\end{define}\newpage
\printbibliography[
  heading=bibintoc,   
  title={References}     % change “References” text if needed
]
% `heading=bibnumbered` if you prefer it numbered chapter (in \tableofcontents)
% `heading=bibintoc`    if you prefer it unnumbered but still in the ToC\newpage








\begin{enumerate} 
  \item Rewrite each of these statements in the form \inlinequote{if $p$, then $q$}
  in English.
  \begin{enumerate}
    \item It is necessary to wash the boss's car to get promoted.
    \begin{subanswer}
      If you get promoted, then you wash the boss's car.
    \end{subanswer}
    \item A sufficient condition for enrollment is that you pay the fees.
    \begin{subanswer}
      If you pay the fees, then you are enrolled. (Or alternatively, \inlinequote{can enroll} in the context of lanaguage.)
    \end{subanswer}
    \item You can access the library website only if you are a student.
    \begin{subanswer}
      If you can access the library website, then you are a student.
    \end{subanswer}
    \item The beach erodes whenever there is a storm.
    \begin{subanswer}
      If there is a storm, then the beach erodes.
    \end{subanswer}
  \end{enumerate}

  \newpage


  \item Let $p$ and $q$ be the propositions $p$: \inlinequote{You buy an ice cream cone.}
  and $q$: \inlinequote{It is hot outside.}\\
  {\color{redcolor}\textbf{Note:} For clarity and improved readability, 
  the logical negation of a proposition $p$ may be represented 
  interchangeably as $\sNOT{p}$ or $\NOT p$.}
  \begin{enumerate}[label=(\alph*)]
    \item Write the converse, inverse, and contrapositive of the conditional
    $p \THEN q$ in English using the \inlinequote{if\ldots, then\ldots} form.  
    Clearly label each one.
    \begin{subanswer}
      The converse, inverse, and contrapositive of the conditional $p \THEN q$ in English are listed as follows.
      \begin{itemize}
        \item \textbf{Conditional ($p \THEN q$):} If you buy an ice cream cone, then it is hot outside.
        \item \textbf{Converse ($q \THEN p$):} If it is hot outside, then you buy an ice cream cone.
        \item \textbf{Inverse ($\sNOT{p} \THEN \sNOT{q}$):} If you do not buy an ice cream cone, then it is not hot outside.
        \item \textbf{Contrapositive ($\sNOT{q} \THEN \sNOT{p}$):} If it is not hot outside, then you do not buy an ice cream cone.
      \end{itemize}
    \end{subanswer}
    \item Which, if any,
    of the four implications are logically equivalent to each other?
    \begin{subanswer}
      The original conditional ($p \THEN q$) is logically equivalent to its 
      contrapositive ($\sNOT{q} \THEN \sNOT{p} $). The converse ($q \THEN p$) 
      is logically equivalent to the inverse ($\sNOT{p} \THEN \sNOT{q}$).

    \end{subanswer}
    \item Write the negation of $p \THEN q$ in English as simply as possible, i.e.\ 
    without using \inlinequote{It is not the case that\ldots}.
    \begin{subanswer}
      You buy an ice cream cone and it is not hot outside.
    \end{subanswer}
  \end{enumerate}

  \newpage

















  \item Find a compound proposition $Q$ in terms of propositions $p$, $q$, 
  and $r$ and logical connectives ($\NOT$, $\AND$, $\OR$, $\THEN$, $\IFF$)
  satisfying the conditions.
  Justify your answers with a truth table.\\
  {\color{redcolor}\textbf{Note:} For clarity and improved readability, 
  the logical negation of a proposition $p$ may be represented 
  interchangeably as $\sNOT{p}$ or $\NOT p$.}
  \begin{enumerate}
    \item This $Q$ is true when $p$ and $q$ are true and $r$ is false, 
    but false otherwise.
    \begin{subanswer}
      The compound proposition is $Q \equiv p \AND q \AND \NOT r$. The truth table confirms 
      this as follows.
      \begin{center}
        \begin{tabular}{|c|c|c|c|c|}
        \hline
        $p$ & $q$ & $r$ & $\NOT r$ & $Q$ \\
        \hline
        T & T & T & F & F \\
        T & T & F & T & T \\
        T & F & T & F & F \\
        T & F & F & T & F \\
        F & T & T & F & F \\
        F & T & F & T & F \\
        F & F & T & F & F \\
        F & F & F & T & F \\
        \hline
        \end{tabular}
      \end{center}
    \end{subanswer}
    \item This $Q$ is true when exactly two of $p$, $q$, and $r$ are 
    true, and is false otherwise.
    \begin{subanswer}
      The compound proposition is 
      $Q \equiv (p \AND q \AND \NOT r) \OR (p \AND \NOT q \AND r) \OR (\NOT p \AND q \AND r)$. 
      The truth table confirms the result as follows.
      \begin{center}
        \begin{tabular}{|c|c|c|c|c|c|c|c|c|}
        \hline
        $p$ & $q$ & $r$ & $p \AND q \AND \NOT r$ & $p \AND \NOT q \AND r$ & $\NOT p \AND q \AND r$ & $Q$ \\
        \hline
        T & T & T & F & F & F & F \\
        T & T & F & T & F & F & T \\
        T & F & T & F & T & F & T \\
        T & F & F & F & F & F & F \\
        F & T & T & F & F & T & T \\
        F & T & F & F & F & F & F \\
        F & F & T & F & F & F & F \\
        F & F & F & F & F & F & F \\
        \hline
        \end{tabular}
      \end{center}
    \end{subanswer}
  \end{enumerate}


  \newpage




  \item Let $p$, $q$, and $r$ be propositions.  Determine whether the 
  following pairs of compound propositions are logically equivalent.
  Justify your answer \textbf{without using a full truth table}.\\
  {\color{redcolor}\textbf{Note:} For clarity and improved readability, 
  the logical negation of a proposition $p$ may be represented 
  interchangeably as $\sNOT{p}$ or $\NOT p$.}
  \begin{enumerate}
    \item $(p \THEN q) \AND (p \THEN r)$ and $p \THEN (q \AND r)$.
    \begin{subprove}
      Let \(p,q,r\) be propositions. 
      \[
        (p\THEN q)\AND(p\THEN r) \equiv (\NOT p\OR q)\AND(\NOT p\OR r)
      \]
      by applying the \textbf{Implication Equivalence} \(p\THEN q\equiv~\NOT p\OR q\) to both implications.
      Moreover, by \textbf{Distributive Law} in the form \((p\OR q)\AND(p\OR r)\equiv p\OR(q\AND r)\), 
      \[
        (\NOT p\OR q)\AND(\NOT p\OR r) \equiv \NOT p\OR(q\AND r).
      \]
      Lastly, applying the \textbf{Implication Equivalence} in reverse,
      \[
        \OR{\NOT{p}}{\group{q\AND r}}\equiv\THEN{p}{(q\AND r)}.
      \]
      Thus, \((p\THEN q)\AND(p\THEN r)\equiv p\THEN(q\AND r)\), so the two compound propositions are \textbf{logically equivalent}. 
    \end{subprove}

    \item $(p \THEN r) \AND (q \THEN r)$ and $(p \OR q) \THEN r$.
    \begin{subprove}
      Let \(p,q,r\) be propositions. 
      \[
        (p\THEN r)\AND(q\THEN r) \equiv (\NOT p\OR r)\AND(\NOT q\OR r)
      \]
      by applying the \textbf{Implication Equivalence} \(p\THEN q\equiv~\NOT p\OR q\) to both implications.
      Moreover, by \textbf{Distributive Law} in the form \((p\OR q)\AND(p\OR r)\equiv p\OR(q\AND r)\), 
      considering \textbf{Commutative Law},
      \[
        (\NOT p\OR r)\AND(\NOT q\OR r) \equiv (r\OR \NOT p)\AND(r\OR \NOT q) \equiv r\OR(\NOT p\AND \NOT q) \equiv (\NOT p\AND\NOT q)\OR r.
      \]
      Then by \textbf{DeMorgan's Law} \(\sNOT{p\OR q}\equiv\sNOT{p}\AND\sNOT{q}\) in reverse,
      \[
        (\NOT p\AND\NOT q)\OR r\equiv~\sNOT{p\OR q}\OR r.
      \]
      Furthermore, by the \textbf{Implication Equivalence} in reverse,
      \[
        \sNOT{p\OR q}\OR r \equiv (p\OR q)\THEN r.
      \]
      Thus, by transitivity of logical equivalence, so the two compound propositions are \textbf{logically equivalent}.
    \end{subprove}

    \newpage

    \item $(p \THEN q) \THEN r$ and $p \THEN (q \THEN r)$.
    \begin{subprove}
      Let \(p,q,r\) be propositions. For the sake of counterexample, 
      take \(p=\FALSE\), \(q=\TRUE\), and \(r=\FALSE\). Then, 
      \[
        \text{L.H.S: }(\FALSE\THEN\TRUE)\THEN\FALSE \equiv \TRUE\THEN\FALSE \equiv \FALSE,
      \]
      \[
        \text{R.H.S: }\FALSE\THEN(\TRUE\THEN\FALSE) \equiv \FALSE\THEN\FALSE \equiv \TRUE.
      \]
      Since the truth values differ for some propositions arbitrarily chosen (counterexample), the two compound propositions are \textbf{not logically equivalent}.
    \end{subprove}

    \item $(p \OR q) \THEN r$ and $(p \THEN r) \OR (q \THEN r)$.
    \begin{subprove}
      Let \(p,q,r\) be propositions. For the sake of counterexample, 
      take $p=\TRUE$, $q=\FALSE$, and $r=\FALSE$. Then, 
      \[
        \text{L.H.S: }(\TRUE\OR\FALSE)\THEN\FALSE \equiv \TRUE\THEN\FALSE \equiv \FALSE,
      \]
      \[
        \text{R.H.S: }(\TRUE\THEN\FALSE)\OR(\FALSE\THEN\FALSE) \equiv \FALSE\OR\TRUE \equiv \TRUE.
      \]
      Since the truth values differ for some propositions arbitrarily chosen (counterexample), the two compound propositions are \textbf{not logically equivalent}.
    \end{subprove}
  \end{enumerate}






  \newpage





























  \item 
  Prove the following directly from the two definitions and three axioms
  given in ICA2 on 8/20.
  \begin{define}{ICA2 Definitions and Axioms}
    Here are the ICA2 Definitions and Axioms as follows.
    \begin{enumerate}
      \item \textbf{Definition:} An integer $n$ is \underbar{even} if there exists some integer $k$ such that
      \[
      n = 2k.
      \]
      \item \textbf{Definition:} An integer $n$ is \underbar{odd} if there exists some integer $k$ such that
      \[
      n = 2k + 1.
      \]
      \item \textbf{Axiom 1:} The sum of two integers is again an integer. (Closure under addition)
      \item \textbf{Axiom 2:} The product of two integers is again an integer. (Closure under multiplication)
      \item \textbf{Axiom 3:} An integer is odd if and only if it is not even.
    \end{enumerate}
  \end{define}
  You may not use any other results without proving them first.
  Every statement in your proof must be fully justified.
  See \textbf{Elements of (Proof) Style} guidelines in Canvas for information.
  \begin{enumerate}
  \item Direct proof: For all integers $n$, if $n$ is odd, then $9n + 5$ is even.
  \begin{subprove}
    Let $n$ be an arbitrary integer then assume $n$ is odd (condition) for the sake of direct proof. 
    By the definition of odd, there exists an integer $k$ such that
    \[
      n = 2k + 1.
    \]
    By substitution, 
    \[
      9n + 5 = 9(2k + 1) + 5 = 18k + 9 + 5 = 18k + 14= 2(9k + 7).
    \]
    By Axiom 1 and 2 (closure under addition and multiplication), $9k+7$ is an integer since $9$, $k$, and $7$ are integers. (Trivial)

    Therefore, by the definition of even, $9n+5$ is even (conclusion).
    Hence, by direct proof, \inlinequote{for all integers $n$, if $n$ is odd, then $9n + 5$ is even.}
  \end{subprove}



  \item Contrapositive proof: For all integers $n$, if $7n + 5$ is odd, then $n$ is even.
  \begin{subprove}
    Let $n$ be an arbitrary integer then assume $n$ is odd (\NOT conclusion) for the sake of contrapositive proof.
    By the definition of odd, there exists an integer $k$ such that
    \[
      n = 2k + 1.
    \]
    By substitution,
    \[
      7n + 5 = 7(2k + 1) + 5 = 14k + 7 + 5 = 14k + 12 = 2(7k + 6).
    \]
    By Axiom 1 and Axiom 2 (closure under addition and multiplication), 
    $7k+6$ is an integer since $7$, $k$, and $6$ are integers. (Trivial)

    Therefore, by the definition of even, $7n+5$ is even (\NOT condition).
    Hence, by contrapositive proof, \inlinequote{for all integers $n$, if $7n + 5$ is odd, then $n$ is even.}
  \end{subprove}

  \item Proof by contradiction: For all integers $n$, if $3n + 2$ is even, then $n$ is even.
  \begin{subprove}
    Let $n$ be an arbitrary integer then assume $3n+2$ is even and assume, for the sake of contradiction, that $n$ is not even.
    By Axiom 3, $n$ not even implies $n$ is odd. By the definition of odd, there exists an integer $k$ such that
    \[
      n = 2k + 1.
    \]
    By substitution,
    \[
      3n + 2 = 3(2k + 1) + 2 = 6k + 3 + 2 = 6k + 5= 2(3k+2) + 1.
    \]
    By Axiom 1 and Axiom 2 (closure under addition and multiplication), 
    $3k+2$ is an integer since $3$, $k$, and $2$ are integers. (Trivial)
    
    Thus $3k+2$ is odd by the definition of odd, and this contradicts the assumption that 
    $3n+2$ is even. 
    Hence, by contradiction, \inlinequote{for all integers $n$, if $3n + 2$ is even, then $n$ is even.}
  \end{subprove}
  \end{enumerate}



  \newpage



  \item
  Meanwhile, back in the Land of Logicalia,
  before you are allowed into the castle, you must identify the
  type of the two guards: Branwyn (she) and Caradoc (he). \\
  \textbf{Prove} your answers to the puzzles are correct.
  Be certain to first clearly state your claim.

  \begin{enumerate}
  \item
  Branwyn says \inlinequote{If I am a Knight, then so is Caradoc.}
  Can you enter the castle?\\

  {\color{redcolor}\textbf{Claim}: \inlinequote{Branwyn is a Knight and Caradoc is a Knight.}}
  \begin{subprove}
    Let the propositions \(p\) be \inlinequote{Branwyn is a Knight} and 
    \(q\) be \inlinequote{Caradoc is a Knight}. Let \(S\) be Branwyn's propositional 
    sentence, the implication given, then it is \(p \THEN q\).\\
    \textbf{Inferences of Implications}
    \begin{enumerate}
      \item Assume, for contradiction, \(\NOT p\). (Branwyn is a Knave.)
      \item \(S\) is false. (Knave always lies.)
      \item \(p \THEN q\) is false, then, \(p\) is true and \(q\) is false.
      \item \(\NOT p\) and \(p\), so \FALSE. Therefore, the assumption \(\NOT p\) is impossible.
      \item Thus, \(p\) is true. (by contradiction)
      \item If \(p\) is true, then \(S\) was made by a Knight, so \(S\) is true. (Knights never lies.)
      \item The conditional, \(S\), is true and the condition, \(p\), is true. Thus, the conclusion, $q$, is true.
      \item Therefore, \(p\) and \(q\).
    \end{enumerate}
    Hence, Branwyn is a Knight and Caradoc is a Knight. (Yes, can enter)
  \end{subprove}


  \item What if Caradoc had instead said \inlinequote{If Branwyn is a Knight, then I am a Knave.}
  Could you enter the castle now?\\

  {\color{redcolor}\textbf{Claim}: Branwyn is a Knave and Caradoc is a Knight.}
  \begin{subprove}
    Let the propositions \(p\) be \inlinequote{Branwyn is a Knight} and 
    \(q\) be \inlinequote{Caradoc is a Knight}. Let \(S\) be Caradoc's propositional
    sentence, the implication given, then it is \(p \THEN \NOT q\).\\
    \textbf{Inferences of Implications}
    \begin{enumerate}
      \item Assume, for contradiction, \(\NOT q\). (Caradoc is a Knave.)
      \item \(S\) is false. (Knave always lies.)
      \item \(p \THEN ~\NOT q\) is false, thus, \(p\) is true and \(q\) is true.
      \item \(\NOT q\) and \(q\), so \FALSE. Therefore, the assumption \(\NOT q\) is impossible.
      \item Thus, \(q\) is true. (by contradiction)
      \item If \(q\) is true, then \(S\) was made by a Knight, so \(S\) is true. (Knights never lie.)
      \item The conditional, \(S\), is true. Moreover, its contrapositive is also true. So, \(q \THEN \NOT p\) is true.
      \item Since \(q \THEN \NOT p\) and its condition, $q$, is true, its conclusion, \(\NOT p\), is also true.
      \item Therefore, \(\NOT p\) and \(q\).
    \end{enumerate}
    Hence, Branwyn is a Knave and Caradoc is a Knight. (No, cannot enter)
  \end{subprove}
  \end{enumerate}




  \newpage

  The following definitions and axioms are from ICA2 on 8/20.

  \begin{define}{ICA2 Definitions and Axioms}
    Here are the ICA2 Definitions and Axioms as follows.
    \begin{enumerate}
      \item \textbf{Definition:} An integer $n$ is \underbar{even} if 
      there exists some integer $k$ such that
      \[
      n = 2k.
      \]
      \item \textbf{Definition:} An integer $n$ is \underbar{odd} if 
      there exists some integer $k$ such that
      \[
      n = 2k + 1.
      \]
      \item \textbf{Axiom 1:} The sum of two integers is again an 
      integer. (Closure under addition)
      \item \textbf{Axiom 2:} The product of two integers is again 
      an integer. (Closure under multiplication)
      \item \textbf{Axiom 3:} An integer is odd if and only if it 
      is not even.
    \end{enumerate}
  \end{define}

  Moreover, the following definitions are provided from in-class 
  activities (ICA) and in-webwork contents.

  \begin{define}{Divisibility}
    An integer $n$ divides an integer $m$, denoted $n \mid m$, if there 
    exist an integer $k$ such that $m=n k$. Otherwise, $n \nmid m$.
  \end{define}

  \begin{define}{Set Difference}
    Let $A$, $B$ be subsets of a universal set $\mathcal{U}$.
    The set difference $A \setminus B$ is defined to be the set
    $\{ x \in \mathcal{U} \mid x \in A \text{ and } x \notin B\}$. \\
    \textit{It is sometimes called the relative complement of $B$ in $A$, and 
    denoted in the textbook as $A - B$.}
  \end{define}

  \item Let $a \in \Z$.  Prove that $3 \mid a^2$ if and only if $3 \mid a$.
  \label{lemma1}
  \textit{You may use the fact that every integer can be written as 
  exactly one of $3k$, $3k + 1$, or $3k+2$ for some integer $k$.}
  \begin{subprove}
    Let $a \in \Z$ be arbitrary. We will prove both directions of the biconditional.
    
    \textbf{Forward direction:} Assume $3 \mid a^2$. We will show $3 \mid a$.
    By the given fact, $a$ can be written as exactly one of $3k$, $3k + 1$, or $3k+2$ for some integer $k$.
    \begin{itemize}
      \item If $a = 3k$, then $3 \mid a$ by definition.
      \item If $a = 3k + 1$, then $a^2 = (3k + 1)^2 = 9k^2 + 6k + 1 = 3(3k^2 + 2k) + 1$. Since $3k^2 + 2k$ is an integer, $a^2 = 3m + 1$ for some integer $m$, so $3 \nmid a^2$, contradicting our assumption.
      \item If $a = 3k + 2$, then $a^2 = (3k + 2)^2 = 9k^2 + 12k + 4 = 3(3k^2 + 4k + 1) + 1$. Since $3k^2 + 4k + 1$ is an integer, $a^2 = 3m + 1$ for some integer $m$, so $3 \nmid a^2$, contradicting our assumption.
    \end{itemize}
    Therefore, $a$ must be of the form $3k$, so $3 \mid a$.
    
    \textbf{Reverse direction:} Assume $3 \mid a$. We will show $3 \mid a^2$.
    Since $3 \mid a$, there exists an integer $k$ such that $a = 3k$.
    Then $a^2 = (3k)^2 = 9k^2 = 3(3k^2)$.
    Since $3k^2$ is an integer, $3 \mid a^2$.
    
    Hence, $3 \mid a^2$ if and only if $3 \mid a$.
  \end{subprove}

  \item Prove that $\sqrt{3}$ is irrational.  
  \textit{Hint: use Problem~\ref{lemma1}.}
  \begin{subprove}
    Assume, for the sake of contradiction, that $\sqrt{3}$ is rational.
    Then there exist integers $p$ and $q$ with $q \neq 0$ such that $\sqrt{3} = \frac{p}{q}$.
    Without loss of generality, we may assume that $p$ and $q$ are coprime (have no common factors other than $\pm 1$).
    
    Squaring both sides, we get $3 = \frac{p^2}{q^2}$, which implies $3q^2 = p^2$.
    This means $3 \mid p^2$.
    
    By Problem~\ref{lemma1}, since $3 \mid p^2$, we have $3 \mid p$.
    Therefore, there exists an integer $k$ such that $p = 3k$.
    
    Substituting back, we get $3q^2 = (3k)^2 = 9k^2$, which simplifies to $q^2 = 3k^2$.
    This means $3 \mid q^2$.
    
    Again by Problem~\ref{lemma1}, since $3 \mid q^2$, we have $3 \mid q$.
    Therefore, there exists an integer $m$ such that $q = 3m$.
    
    But this contradicts our assumption that $p$ and $q$ are coprime, since both are divisible by $3$.
    Hence, our assumption that $\sqrt{3}$ is rational must be false.
    
    Therefore, $\sqrt{3}$ is irrational.
  \end{subprove}

  \item Prove the following directly from the definition.   \\
  Let $a$, $b$, $c$, $d$, $x$ and $y$ be integers with $a \neq 0$ and $b \neq 0$.
  \begin{enumerate}
    \item If $a \mid c$ and $b \mid d$, then $ab \mid cd$.
    \begin{subprove}
      Let $a$, $b$, $c$, $d$ be integers with $a \neq 0$ and $b \neq 0$.
      Assume $a \mid c$ and $b \mid d$.
      By definition of divisibility, there exist integers $k$ and $l$ such that $c = ak$ and $d = bl$.
      Then $cd = (ak)(bl) = (ab)(kl)$.
      Since $kl$ is an integer, by definition of divisibility, $ab \mid cd$.
    \end{subprove}
    
    \item If $a \mid c$ and $a \mid d$, then $a \mid cx + dy$.
    \begin{subprove}
      Let $a$, $c$, $d$, $x$, $y$ be integers with $a \neq 0$.
      Assume $a \mid c$ and $a \mid d$.
      By definition of divisibility, there exist integers $k$ and $l$ such that $c = ak$ and $d = al$.
      Then $cx + dy = (ak)x + (al)y = a(kx + ly)$.
      Since $kx + ly$ is an integer, by definition of divisibility, $a \mid cx + dy$.
    \end{subprove}
    
    \item If $a \mid b$ and $b \mid a$, then $a = b$ or $a = -b$.
    \begin{subprove}
      Let $a$, $b$ be integers with $a \neq 0$ and $b \neq 0$.
      Assume $a \mid b$ and $b \mid a$.
      By definition of divisibility, there exist integers $k$ and $l$ such that $b = ak$ and $a = bl$.
      Substituting the second equation into the first, we get $b = (bl)k = b(lk)$.
      Since $b \neq 0$, we can divide both sides by $b$ to get $1 = lk$.
      The only integer solutions to this equation are $(l,k) = (1,1)$ or $(l,k) = (-1,-1)$.
      If $(l,k) = (1,1)$, then $a = b$.
      If $(l,k) = (-1,-1)$, then $a = -b$.
      Therefore, $a = b$ or $a = -b$.
    \end{subprove}
    
    \item If $a \nmid cd$, then $a \nmid c$ and $a \nmid d$.
    \begin{subprove}
      Let $a$, $c$, $d$ be integers with $a \neq 0$.
      We will prove the contrapositive: if $a \mid c$ or $a \mid d$, then $a \mid cd$.
      Assume $a \mid c$ or $a \mid d$.
      \begin{itemize}
        \item If $a \mid c$, then there exists an integer $k$ such that $c = ak$. Then $cd = (ak)d = a(kd)$. Since $kd$ is an integer, $a \mid cd$.
        \item If $a \mid d$, then there exists an integer $l$ such that $d = al$. Then $cd = c(al) = a(cl)$. Since $cl$ is an integer, $a \mid cd$.
      \end{itemize}
      In both cases, $a \mid cd$.
      Therefore, by contrapositive, if $a \nmid cd$, then $a \nmid c$ and $a \nmid d$.
    \end{subprove}
  \end{enumerate}

  \item Let $A$, $B$ be subsets of a universal set $\mathcal{U}$.
  Prove the following using the element method:
  \begin{enumerate}
    \item $A \setminus B \subseteq A$.
    \begin{subprove}
      Let $A$, $B$ be subsets of a universal set $\mathcal{U}$.
      To show $A \setminus B \subseteq A$, we need to show that every element in $A \setminus B$ is also in $A$.
      Let $x \in A \setminus B$ be arbitrary.
      By definition of set difference, $x \in A$ and $x \notin B$.
      Since $x \in A$, we have shown that $A \setminus B \subseteq A$.
    \end{subprove}
    
    \item $A \setminus B = A \cap \sNOT{B}$
    \begin{subprove}
      Let $A$, $B$ be subsets of a universal set $\mathcal{U}$.
      We will show that $A \setminus B = A \cap \sNOT{B}$ by showing both containments.
      
      \textbf{Forward containment:} Let $x \in A \setminus B$ be arbitrary.
      By definition of set difference, $x \in A$ and $x \notin B$.
      Since $x \notin B$, we have $x \in \sNOT{B}$.
      Therefore, $x \in A$ and $x \in \sNOT{B}$, so $x \in A \cap \sNOT{B}$.
      Hence, $A \setminus B \subseteq A \cap \sNOT{B}$.
      
      \textbf{Reverse containment:} Let $x \in A \cap \sNOT{B}$ be arbitrary.
      By definition of intersection, $x \in A$ and $x \in \sNOT{B}$.
      Since $x \in \sNOT{B}$, we have $x \notin B$.
      Therefore, $x \in A$ and $x \notin B$, so $x \in A \setminus B$.
      Hence, $A \cap \sNOT{B} \subseteq A \setminus B$.
      
      Since both containments hold, $A \setminus B = A \cap \sNOT{B}$.
    \end{subprove}
    
    \item $A \cap (B \setminus A) = \emptyset$.
    \begin{subprove}
      Let $A$, $B$ be subsets of a universal set $\mathcal{U}$.
      We will show that $A \cap (B \setminus A) = \emptyset$ by contradiction.
      Assume, for the sake of contradiction, that there exists an element $x \in A \cap (B \setminus A)$.
      Then $x \in A$ and $x \in B \setminus A$.
      Since $x \in B \setminus A$, we have $x \in B$ and $x \notin A$.
      But this gives us $x \in A$ and $x \notin A$, which is a contradiction.
      Therefore, our assumption must be false, and $A \cap (B \setminus A) = \emptyset$.
    \end{subprove}
    
    \item $A \cup (B \setminus A) = A \cup B$.
    \begin{subprove}
      Let $A$, $B$ be subsets of a universal set $\mathcal{U}$.
      We will show that $A \cup (B \setminus A) = A \cup B$ by showing both containments.
      
      \textbf{Forward containment:} Let $x \in A \cup (B \setminus A)$ be arbitrary.
      Then $x \in A$ or $x \in B \setminus A$.
      If $x \in A$, then $x \in A \cup B$.
      If $x \in B \setminus A$, then $x \in B$ and $x \notin A$, so $x \in B$, which means $x \in A \cup B$.
      In both cases, $x \in A \cup B$. Hence, $A \cup (B \setminus A) \subseteq A \cup B$.
      
      \textbf{Reverse containment:} Let $x \in A \cup B$ be arbitrary.
      Then $x \in A$ or $x \in B$.
      If $x \in A$, then $x \in A \cup (B \setminus A)$.
      If $x \in B$, then either $x \in A$ or $x \notin A$.
      If $x \in A$, then $x \in A \cup (B \setminus A)$.
      If $x \notin A$, then $x \in B \setminus A$, so $x \in A \cup (B \setminus A)$.
      In all cases, $x \in A \cup (B \setminus A)$. Hence, $A \cup B \subseteq A \cup (B \setminus A)$.
      
      Since both containments hold, $A \cup (B \setminus A) = A \cup B$.
    \end{subprove}
  \end{enumerate}

  \item Let $A$, $B$, and $C$ be subsets of a universal set $\mathcal{U}$.
  Prove or disprove: 
  \begin{enumerate}
    \item If \(A \cap B = A \cap C\), then \(B = C\).
    \begin{subprove}
      This statement is \textbf{false}. We will provide a counterexample.
      Let $\mathcal{U} = \{1, 2, 3, 4\}$, $A = \{1, 2\}$, $B = \{1, 3\}$, and $C = \{1, 4\}$.
      Then $A \cap B = \{1\}$ and $A \cap C = \{1\}$, so $A \cap B = A \cap C$.
      However, $B = \{1, 3\} \neq \{1, 4\} = C$.
      Therefore, the statement is false.
    \end{subprove}
    
    \item If \(A \cup B = A \cup C\), then \(B = C\).
    \begin{subprove}
      This statement is \textbf{false}. We will provide a counterexample.
      Let $\mathcal{U} = \{1, 2, 3, 4\}$, $A = \{1, 2\}$, $B = \{3\}$, and $C = \{4\}$.
      Then $A \cup B = \{1, 2, 3\}$ and $A \cup C = \{1, 2, 4\}$, so $A \cup B = A \cup C$.
      However, $B = \{3\} \neq \{4\} = C$.
      Therefore, the statement is false.
    \end{subprove}
    
    \item If \(A \cap B = A \cap C\) and \(A \cup B = A \cup C\), then \(B = C\).
    \begin{subprove}
      This statement is \textbf{true}. We will prove it.
      Let $A$, $B$, and $C$ be subsets of a universal set $\mathcal{U}$.
      Assume $A \cap B = A \cap C$ and $A \cup B = A \cup C$.
      
      We will show that $B = C$ by showing both containments.
      
      \textbf{Forward containment:} Let $x \in B$ be arbitrary.
      Since $x \in B$, we have $x \in A \cup B = A \cup C$.
      Therefore, $x \in A$ or $x \in C$.
      \begin{itemize}
        \item If $x \in A$, then $x \in A \cap B = A \cap C$, so $x \in C$.
        \item If $x \in C$, then we are done.
      \end{itemize}
      In both cases, $x \in C$. Hence, $B \subseteq C$.
      
      \textbf{Reverse containment:} Let $x \in C$ be arbitrary.
      Since $x \in C$, we have $x \in A \cup C = A \cup B$.
      Therefore, $x \in A$ or $x \in B$.
      \begin{itemize}
        \item If $x \in A$, then $x \in A \cap C = A \cap B$, so $x \in B$.
        \item If $x \in B$, then we are done.
      \end{itemize}
      In both cases, $x \in B$. Hence, $C \subseteq B$.
      
      Since both containments hold, $B = C$.
    \end{subprove}
  \end{enumerate}

  \item Let $A$, $B$ be subsets of a universal set $\mathcal{U}$.
  Prove or disprove:
  \begin{enumerate}
    \item $A \times B = \emptyset$ if $A = \emptyset$ or $B = \emptyset$.
    \begin{subprove}
      This statement is \textbf{true}. We will prove it.
      Let $A$, $B$ be subsets of a universal set $\mathcal{U}$.
      Assume $A = \emptyset$ or $B = \emptyset$.
      
      \textbf{Case 1:} $A = \emptyset$.
      Then $A \times B = \emptyset \times B = \emptyset$ (since there are no first coordinates).
      
      \textbf{Case 2:} $B = \emptyset$.
      Then $A \times B = A \times \emptyset = \emptyset$ (since there are no second coordinates).
      
      In both cases, $A \times B = \emptyset$.
      Therefore, if $A = \emptyset$ or $B = \emptyset$, then $A \times B = \emptyset$.
    \end{subprove}
    
    \item $A \times B = \emptyset$ only if $A = \emptyset$ or $B = \emptyset$.
    \begin{subprove}
      This statement is \textbf{true}. We will prove it by contrapositive.
      Let $A$, $B$ be subsets of a universal set $\mathcal{U}$.
      We will prove: if $A \neq \emptyset$ and $B \neq \emptyset$, then $A \times B \neq \emptyset$.
      
      Assume $A \neq \emptyset$ and $B \neq \emptyset$.
      Since $A \neq \emptyset$, there exists an element $a \in A$.
      Since $B \neq \emptyset$, there exists an element $b \in B$.
      Then $(a,b) \in A \times B$ by definition of Cartesian product.
      Therefore, $A \times B \neq \emptyset$.
      
      By contrapositive, if $A \times B = \emptyset$, then $A = \emptyset$ or $B = \emptyset$.
    \end{subprove}
  \end{enumerate}

\end{enumerate}